% Part: incompleteness
% Chapter: incompleteness-provability
% Section: lob-thm

\documentclass[../../../include/open-logic-section]{subfiles}

\begin{document}

\olfileid{inc}{inp}{tar}

\olsection{সত্যের অসংজ্ঞেয়তা}

\emph{সংজ্ঞেয়তা}-র ধারণাটি পাটিগণিতের ভাষার একটি বিধিবদ্ধ অর্থতত্ত্বের
উপর নির্ভর করে। আমরা পাটিগণিতের ভাষায় সূত্র ও বাক্যের একটি সেট বর্ণনা
করেছি। ``অভিপ্রেত ব্যাখ্যা'' হলো এই বাক্যগুলিকে স্বাভাবিক সংখ্যা সম্বন্ধে
উক্তি হিসেবে পড়া; এমন উক্তি সত্য বা মিথ্যা হতে পারে। ধরি $\Struct{N}$
হলো !!{structure}, যার সংজ্ঞাক্ষেত্র $\Nat$ এবং পাটিগণিতের ভাষার
প্রতীকগুলির ব্যাখ্যা প্রমিত। তখন $\Sat{N}{!A}$-এর অর্থ ``$!A$ প্রমিত
ব্যাখ্যায় সত্য।''

\begin{defn}
স্বাভাবিক সংখ্যার কোনো সম্বন্ধ $R(x_1,\dots,x_k)$-কে $\Struct{N}$-এ
\emph{সংজ্ঞেয়} বলা হয় যদি এবং কেবল যদি পাটিগণিতের ভাষায় এমন একটি
সূত্র $!A(x_1,\dots,x_k)$ থাকে যাতে প্রতিটি $n_1,\dots,n_k$-এর জন্য
$R(n_1,\dots,n_k)$ তখনই সত্য, যখন $\Sat{N}{!A(\num n_1,\dots,\num
  n_k)}$।
\end{defn}

অন্যভাবে বললে, কোনো সম্বন্ধ $\Struct{N}$-এ সংজ্ঞেয় যদি এবং কেবল যদি
তত্ত্ব $\Th{TA}$-তে সেটি উপস্থাপনযোগ্য; এখানে $\Th{TA} =
\Setabs{!A}{\Sat{N}{!A}}$ হলো পাটিগণিতের সত্য বাক্যগুলির সেট। (এটি
তোমার কাছে অবিলম্বে স্পষ্ট না হলে সংজ্ঞাগুলিতে ফিরে গিয়ে যাচাই করো এবং
নিজেকে বোঝাও যে সত্যিই তা-ই।)

\begin{lem}
প্রতিটি গণনসাধ্য সম্বন্ধ~$\Struct{N}$-এ সংজ্ঞেয়।
\end{lem}

\begin{proof}
সহজেই যাচাই করা যায়, $\Th{Q}$-তে কোনো সম্বন্ধকে উপস্থাপনকারী সূত্রটি
$\Struct{N}$-এ একই সম্বন্ধকে সংজ্ঞায়িত করে।
\end{proof}

এবার বিপরীতটিও সত্য কি না জিজ্ঞাসা করা যায়। অর্থাৎ $\Struct{N}$-এ
সংজ্ঞেয় প্রতিটি সম্বন্ধ কি গণনসাধ্য? উত্তর না। যেমন:

\begin{lem}
হল্টিং সম্বন্ধটি $\Struct{N}$-এ সংজ্ঞেয়।
\end{lem}

\begin{proof}
মনে রেখো, ক্লিনি স্বাভাবিক রূপ উপপাদ্য বলে, প্রতিটি আংশিক গণনসাধ্য
অপেক্ষক~$f$-এর এমন একটি সূচক~$e$ আছে যাতে প্রতিটি $x \in \Nat$-এর জন্য
$f(x) = U(\umin{s}{T(e,x,s)})$; এখানে $U$ ও $T$ আদিম পুনরাবৃত্ত এবং তাই
পূর্ণ। অতএব $f(x)$ তখনই সংজ্ঞায়িত (অর্থাৎ গণনা তখনই থামে), যখন এমন
কোনো~$s$ আছে যাতে $T(e,x,s)$ সত্য।

এবার $H$ হলো হল্টিং সম্বন্ধ, অর্থাৎ
\[
H = \Setabs{\tuple{e,x}}{\lexists[s][T(e, x, s)]}.
\]
ধরি $!D_T$ সূত্রটি $\Struct{N}$-এ $T$-কে সংজ্ঞায়িত করে। তখন
\[
H = \Setabs{\tuple{e,x}}{\Sat{N}{\lexists[s][!D_T(\num e, \num x, s)]}},
\]
তাই $\lexists[s][!D_T(z, x, s)]$ সূত্রটি $\Struct{N}$-এ~$H$-কে
সংজ্ঞায়িত করে।
\end{proof}

\begin{prob}
দেখাও যে $Q(n) \defiff n \in \Setabs{\Gn{!A}}{\Th{Q} \Proves !A}$
পাটিগণিতে সংজ্ঞেয়।
\end{prob}

$\Th{TA}$ নিজে সম্বন্ধে কী বলা যায়? সেটি কি পাটিগণিতে সংজ্ঞেয়? অর্থাৎ
সেট $\Setabs{\Gn{!A}}{\Sat{N}{!A}}$ কি পাটিগণিতে সংজ্ঞেয়? তারস্কির
উপপাদ্য এর নেতিবাচক উত্তর দেয়।

\begin{thm}
\ollabel{thm:tarski}
পাটিগণিতের সত্য !!{sentence}s-এর সেটটি পাটিগণিতে সংজ্ঞেয় নয়।
\end{thm}

\begin{proof}
ধরি $!D(x)$ সেটিকে সংজ্ঞায়িত করে; অর্থাৎ $\Sat{N}{!A}$ তখনই সত্য, যখন
$\Sat{N}{!D(\gn{!A})}$। স্থির-বিন্দু লেমা অনুসারে এমন একটি সূত্র $!A$
আছে যাতে $\Th{Q} \Proves !A \liff \lnot !D(\gn{!A})$, এবং তাই
$\Sat{N}{!A \liff \lnot !D(\gn{!A})}$। কিন্তু তখন $\Sat{N}{!A}$ যদি
এবং কেবল যদি $\Sat{N}{\lnot !D(\gn{!A})}$; অথচ $!D(y)$-এর পাটিগণিতের
সত্য উক্তিগুলির সেটকে সংজ্ঞায়িত করার কথা—এর সঙ্গে বিরোধ হয়।
\end{proof}

তারস্কি এই বিশ্লেষণটি সত্যের আরও সাধারণ দার্শনিক ধারণায় প্রয়োগ করেন।
যেকোনো ভাষা $L$ দেওয়া থাকলে তারস্কি যুক্তি দেন, $L$-এর জন্য সত্যের
পর্যাপ্ত ধারণাকে প্রতিটি বাক্য $X$-এর জন্য নিচেরটি পূরণ করতে হবে:
\begin{quote}
`$X$' সত্য যদি এবং কেবল যদি $X$।
\end{quote}
ইংরেজির জন্য তারস্কির বহুল উদ্ধৃত উদাহরণটির বাংলা রূপ হলো:
\begin{quote}
`তুষার সাদা' সত্য যদি এবং কেবল যদি তুষার সাদা।
\end{quote}
কিন্তু কর্ণীকরণ অপেক্ষককে উপস্থাপন করার মতো যথেষ্ট শক্তিশালী যেকোনো ভাষা এবং
যেকোনো ভাষাগত বিধেয় $T(x)$-এর জন্য আমরা ``$X$ যদি এবং কেবল যদি
$T(\text{`$X$'})$ না হয়'' পূরণকারী একটি বাক্য $X$ নির্মাণ করতে পারি।
কোনো সত্যতা-বিধেয় কোনো বাক্যকে একই সঙ্গে সত্য ও মিথ্যা ঘোষণা করুক—এটি
আমরা চাই না; তাই তারস্কির সিদ্ধান্ত ছিল, কোনো ভাষার সীমানার বাইরে কোনোভাবে
না গিয়ে ওই ভাষার সব বাক্যের জন্য সত্যতা-বিধেয় নির্দিষ্ট করা যায় না।
অন্য কথায়, কোনো ভাষার সত্যতা-বিধেয় সেই ভাষাতেই সংজ্ঞায়িত করা যায় না।

\end{document}
